\chapter{Softwarová architektura systému a model konkurence}
\label{ch:architecture}

Při návrhu a implementaci systému se klade vysoký důraz na kombinaci moderního, responsivního uživatelského rozhraní a maximálního výpočetního výkonu výpočetního jádra. Architektura je proto koncipována jako distribuovaná, kde je prezentační vrstva striktně oddělena od nízkoúrovňových matematických a souborových operací. Tento přístup umožňuje efektivně zpracovávat gigabajtové sady astronomických dat bez negativního dopadu na plynulost odezvy aplikace.

\section{Využití frameworku Tauri a jazyka Rust} Architektura systému je navržena tak, aby efektivně propojovala flexibilitu moderních webových technologií s nekompromisním výpočetním výkonem nízkoúrovňového systémového programování. Jako základní stavební kámen klientské prezentační části (frontend) byl zvolen reaktivní ekosystém SvelteKit\footnote{\url{https://svelte.dev/}}. Tento framework vyniká vysokou rychlostí kompilace, efektivním vykreslováním bez nutnosti použití virtuálního \textbf{\ac{DOM}} a vysoce přehlednou správou stavu uživatelského rozhraní. Využití webových standardů pro klientskou část přináší aplikaci nezbytnou modularitu a plynulost při tvorbě responsivního designu i komplexních interaktivních vizualizačních panelů.
Aby však aplikace dokázala bez latence manipulovat s rozsáhlými vícedimenzionálními maticemi astronomických dat, bylo absolutně nezbytné fyzicky oddělit prezentační vrstvu od výpočetně náročných operací. Veškeré algoritmy pro čtení \ac{FITS} souborů, pixelovou aritmetiku, statistické vyhodnocování a extrakci metrik jsou proto plně implementovány v moderním kompilovaném jazyce Rust\footnote{\url{https://www.rust-lang.org/}}. Tento jazyk poskytuje exaktní kontrolu nad alokací paměti a výpočetní rychlost plně srovnatelnou s jazyky C/C++. Jeho striktní kompilátor navíc matematicky garantuje absolutní paměťovou a vláknovou bezpečnost (\textit{memory and thread safety}) bez nutnosti omezujícího běhu správce paměti (\textit{Garbage Collector}). Jak dokládá interní architektura aplikace, Rust umožňuje masivní paralelizaci výpočtů (pomocí vláknového fondu knihovny \textit{Rayon}) a paměťově úsporné zpracování (využívající datové proudy a efektivní lineární algoritmy), což je pro redukci velkých astrofotografických datových sad klíčové.
Kombinaci a bezpečné propojení těchto dvou zcela odlišných prostředí zajišťuje moderní framework Tauri\footnote{\url{https://tauri.app/}}. Na rozdíl od starších, paměťově neefektivních platforem (jako je například Electron), které s sebou přenášejí kompletní instanci prohlížeče Chromium a masivní běžící proces Node.js, Tauri využívá odlehčené nativní vykreslovací jádro hostitelského operačního systému (\textit{WebView}). Komunikace mezi asynchronním frontendem postaveným na SvelteKitu a vysoce výkonným backendem v Rustu probíhá formou serializovaných zpráv přes zabezpečené meziprocesové rozhraní (\textbf{\ac{IPC}}). Toto hybridní spojení umožňuje plně využít pokročilé návrhové vzory webového vývoje a současně dosáhnout maximální hardwarové propustnosti při zpracování obrazových dat na úrovni operačního systému.

\section{Paralelizace pomocí knihovny Rayon}
Surové snímky ve formátu \ac{FITS} obsahují miliony pixelů, nad kterými je nutné v rámci kalibrace, mřížkového odhadu pozadí a výpočtu statistik provádět opakované maticové transformace. Sekvenční zpracování takového objemu dat na jednom procesorovém jádru by představovalo kritický limit propustnosti celého systému.

Za účelem maximálního využití moderního vícejádrového hardwaru je výpočetní jádro paralelizováno pomocí knihovny Rayon\footnote{\url{https://crates.io/crates/rayon}}. Tato knihovna implementuje datový paralelismus založený na pokročilém plánovači krádeže práce (\textit{work-stealing scheduler}). V praxi to umožňuje transformovat standardní sekvenční iterátory nad pixelovými poli a souborovými frontami na konkurentní operace pouhou záměnou metod. 

Algoritmus automaticky a dynamicky dělí pole pixelů na optimální sub-matice, které distribuuje napříč všemi dostupnými logickými jádry CPU. Vzhledem k tomu, že jazyk Rust na úrovni kompilace garantuje absenci souběhů nad daty (\textit{data races}), probíhá toto masivní paralelní zpracování s minimální režií na zamykání vláken. Výpočetní výkon tak lineárně škáluje s počtem procesorových jader hostitelského počítače, což zkracuje proces analýzy a kalibrace desítek snímků z minut na řád jednotek sekund.